\section{光前波函数、介子形状因子、DY 与 Sivers 比}\label{sec:yzliu-exclusive}

同一个重—轻核 $H$ 与同一个 $S_I$ 出现在一系列格点可计算的对象中。本节四个命题的结构相同：领头区域分析给出因子化，$\tilde K$ 与 $S_I$ 在适当的组合中相消。

\begin{proposition}[准光前波函数振幅的因子化 \tagA]\label{prop:yzliu-lfwf}
在微扰论各阶（硬核以部分子态计算）：(i) 以广义软函数 $S^\pm_N$ 定义的、快度重整化的 $N$ 部分子光前波函数振幅 $\psi^\pm_N(x_i,\vec b_{i\perp},\mu,\zeta)$（$\pm$：类光规范连线指向未来／过去）满足 $2\zeta\partial_\zeta\ln\psi^\pm_N=\tilde K_N$ 与一个 RGE；未来与过去两种版本由解析延拓相联，例如 $S^{-\dagger}_N(\delta^+\delta^-)=S^+_N(-\delta^+\delta^-+\ii0)$。(ii) 准光前波函数振幅（类空订书钉）满足
\begin{equation}\label{eq:yzliu-lfwf}
\tilde\psi^\pm_N\sqrt{S_{rN}}=\exp\Big[\frac12\ln\frac{\mp\zeta_z+\ii0}{\zeta}\,\tilde K_N\Big]H^\pm_N\Big(\frac{\zeta_{z,i}}{\mu^2}\Big)\psi^\pm_N+\order{\frac{\Lambda^2}{\zeta_{z,i}},\frac{M^2}{\zeta_{z,i}},\frac1{\delta b_{ij\perp}^2\zeta_{z,i}}},
\end{equation}
对 $q\bar q$ 介子，单圈快度对数相差 $\frac12\tilde K\ln(\zeta_z/\zeta)$（因 $\sqrt{\zeta_z\bar\zeta_z/(\zeta\bar\zeta)}=\zeta_z/\zeta$），$H^\pm_{q\bar q}=C^\pm_{HL}(\zeta_z)\,C^\pm_{HL}(\bar\zeta_z)$（\cref{prop:yzliu-hard-kernel}(b)），其虚部不相消。(iii) $\tilde K_N$ 可由两个动量下的准光前波函数之比得到。非微扰地位：\cref{asm:intro-bridge}、\cref{asm:intro-composite}（MA0、MA1、MA2）。
% errtype: E-S
\uses{thm:lamet-quasi-tmd,prop:yzliu-hard-kernel,thm:yzliu-time-ordering}
\end{proposition}

\begin{proofsketch}
文献\cite[第~II--V~节]{Ji:2021znw}给出证明概要。领头区域与\cref{thm:lamet-quasi-tmd}相同：硬核局限在两个夸克—连线顶点附近，共线子图与光前振幅相同，软辐射由 $S^\pm_N$ 吸收；$\tilde K_N$ 与 TMD 的 $\tilde K$ 相同（$N=2$）是\cref{thm:yzliu-time-ordering}所支持的普适性。$\pm\ii0$ 由光前无穷远处的指数衰减确定。单圈检验：把期刊版的单圈振幅、准振幅、约化软函数与核代入\cref{eq:yzliu-lfwf}，$L_b$ 依赖相消当且仅当指数系数为 $\frac12$ 且取 $+\ii0$；对 $\psi^+$ 取 $-\ii0$ 会留下与 $b_\perp$ 有关、不能被 $H^+$ 吸收的余项 $-2\pi\ii\,aL_b$。
\end{proofsketch}

\paragraph{地位与缺口。}全阶论证 [A]，结构与\cref{thm:lamet-quasi-tmd}相同，最弱的一步仍是幂次计数（E-S）；$\tilde K$ 在光前波函数与 TMD 之间的普适性由\cref{thm:yzliu-time-ordering}支持。指数中的 $\frac12$ 见期刊版\cite[式~(49)、(51)、(74)]{Ji:2021znw}，其脚注~1 指出了其他文献中的遗漏（附录~Q 第~11~条）；期刊版式~(49)、(74) 中的 $-\ii0$ 对 $\psi^+$ 与其单圈结果不相容，$+\ii0$ 的写法见\cite[式~(4.4)]{Liu:2023onm}。与准 TMD 不同，$H^\pm_{q\bar q}$ 的两个因子相位相同（\cref{ex:yzliu-hl-decomposition}）。概念上，光前量子化可看作积掉小 $k^+$ 模式的有效理论，LaMET 则是“有效光前量子化”\cite[第~II.C~节]{Ji:2021znw}。
% Verification: claim E02 (journal version correct; +i0 consistent with the one-loop expressions); script s07 (W26).

\begin{proposition}[大动量下的介子形状因子与 $S_I$ 的提取 \tagA]\label{prop:yzliu-meson-ff}
对由不同味的轻夸克场 $\psi$、$\eta$ 构成的形状因子 $F(b_\perp,P,P',\mu)=\langle P'|\bar\eta(\vec b_\perp)\Gamma'\eta(\vec b_\perp)\,\bar\psi(0)\Gamma\psi(0)|P\rangle$，$P=(P^t,P^z,\vec0_\perp)$、$P'=(P^t,-P^z,\vec0_\perp)$，$\Gamma=\Gamma'=1$ 或 $\gamma_\perp$。在局部假设 (MF) 之下，于 $P^z$ 大、$b_\perp\gg1/P^z$ 时 % VF-YZL-16
\begin{equation}\label{eq:yzliu-meson-ff}
F=S_I(b_\perp,\mu)\int\dd x\,\dd x'\,H(x,x')\,\tilde\psi^\dagger(x',b_\perp)\,\tilde\psi(x,b_\perp)+\cdots,\qquad H=\frac{H_1(Q^2,\bar Q^2,\mu^2)}{H_2(\zeta_z,\bar\zeta_z)\,H_2(\zeta'_z,\bar\zeta'_z)},
\end{equation}
$H_1=H^{\mathrm{Sud}}(-Q^2)H^{\mathrm{Sud}}(-\bar Q^2)$（类空 Sudakov 形状因子的硬核），$Q^2=\sqrt{\zeta_z\zeta'_z}$，$\bar Q^2=\sqrt{\bar\zeta_z\bar\zeta'_z}$，$\tilde\psi$ 是\cref{prop:yzliu-lfwf}的 $q\bar q$ 准光前波函数（订书钉指向 $-\hat z$），$H_2$ 为其匹配核；因此 $S_I=F/\int H\tilde\psi^\dagger\tilde\psi$ 可在格点上计算。(MF)：该遍举形状因子在 Li--Sterman 意义下 TMD 因子化——$0$ 与 $\vec b_\perp$ 处两个硬核、两个共线子图、一个软子图，Glauber 交换如 DY/SIDIS 处理（\cref{lem:tmd-glauber}），全部规范连线取指向未来。非微扰地位：\cref{asm:intro-bridge}、\cref{asm:intro-composite}（MA0、MA1、MA2）。
\uses{prop:yzliu-lfwf,thm:yzliu-intrinsic-soft,lem:tmd-glauber}
\end{proposition}

\begin{proofsketch}
（1）领头区域\cite[第~5~节、图~25]{Liu:2022nnk}：软发散由 $S(b_\perp,\mu,\delta^+,\delta^-)$ 因子化，共线发散由未减除的振幅 $\psi(x,b_\perp,\mu,\delta^{-\prime})$ 吸收并除以软函数以免重复计数；全部连线指向未来时软函数等于标准 DY 软因子，得 $F=\int H_1\,\psi^\dagger\psi\,S/(S\,S)$，正规子全部相消\cite[式~(145)]{Liu:2022nnk}；$H_1$ 由类空 Sudakov 形状因子给出\cite[式~(148)]{Liu:2022nnk}。（2）对准振幅作同样的分析得 $\tilde\psi=H_2\,\psi(\delta^-)S(\delta^+)/S(\delta^+,\delta^-)$\cite[式~(150)]{Liu:2022nnk}。（3）代入并用 $\zeta\zeta'=\zeta_z\zeta'_z$，振幅与 CS 核全部相消，只剩\cref{thm:yzliu-intrinsic-soft}的 $S_I$\cite[式~(151)--(152)]{Liu:2022nnk}。关系 $2\mathcal D_1-\mathcal D_2=\mathcal D$ 与上述主公式已在单圈验证\cite{Ji:2019sxk}。
\end{proofsketch}

\paragraph{地位与缺口。}因子化为全阶论证 [A]；最弱的一步是遍举形状因子的 Glauber 处理，仅类比 DY/SIDIS 论证。2020 年的原始论文\cite{Ji:2019sxk}只陈述了 $F$ 与 $\tilde\psi$ 的因子化而把细节留待以后发表，若只据此，这一关系只能算 [P]；本书依据此后写出的领头区域论证\cite[第~5~节]{Liu:2022nnk}给出 [A]。单圈检验是证据；全部连线指向未来时软函数的时序无关性见\cref{thm:yzliu-time-ordering}。解释：这是 Collins--Soper 轴规范因子化\cite{CollinsSoper1981}的规范不变完成；端点处 $H$ 只有 $1+\alpha_s\ln^2x$ 的对数发散，准波函数线性衰减，端点区行为为 $x\ln^2x$，若令 $x+x'=1$ 则改善为 $x^2\ln^2x$。格点：LPC 在领头阶因子化下计算了 $S_I$\cite{LatticeParton:2020uhz}，单圈修正见\cite{LatticePartonLPC:2023pdv}，单圈 TMD 波函数与软函数见\cite{Deng:2022gzi}；Collins--Soper 核的格点提取更早已有\cite{Shanahan:2020zxr}。

\begin{proposition}[以准 TMD 表示的小 $q_T$ Drell--Yan 截面 \tagA]\label{prop:yzliu-dy-quasi}
在\cref{thm:tmd-factorization}与\cref{thm:lamet-quasi-tmd}成立的条件下，
\begin{equation}\label{eq:yzliu-dy-quasi}
\frac{\dd\sigma_{\mathrm{DY}}}{\dd^2\vec b_\perp\,\dd x\,\dd x'}=\hat\sigma(x,x',\mu,\zeta_z,\zeta'_z)\,S_I(b_\perp,\mu)\,\tilde f(x,b_\perp,\mu,\zeta_z)\,\tilde f(x',b_\perp,\mu,\zeta'_z)+\cdots,
\end{equation}
$\mathcal H(Q^2/\mu^2)=H(\zeta_z/\mu^2)H(\zeta'_z/\mu^2)\hat\sigma$，$\mathcal H$ 为\cref{thm:tmd-factorization}的硬因子；当 $\zeta_z\zeta'_z=\zeta_A\zeta_B=Q^4$ 时 CS 因子相消，不需要取光锥极限。非微扰地位：\cref{asm:intro-bridge}、\cref{asm:intro-composite}（MA0、MA1、MA2）。
\uses{thm:tmd-factorization,thm:lamet-quasi-tmd,thm:yzliu-intrinsic-soft}
\end{proposition}

\begin{proof}
由\cref{thm:tmd-factorization}，$\dd\sigma/(\dd^2\vec b_\perp\dd x\dd x')=\mathcal H\,f^{\mathrm{TMD}}(x,b_\perp,\mu,\zeta_A)f^{\mathrm{TMD}}(x',b_\perp,\mu,\zeta_B)$，$\zeta_A\zeta_B=Q^4$。由\cref{thm:lamet-quasi-tmd}，$f^{\mathrm{TMD}}(\zeta)=\tilde f(\zeta_z)\sqrt{S_r}\,\ee^{-\frac12\tilde K\ln(\zeta_z/\zeta)}/H(\zeta_z/\mu^2)$。两式相乘，指数为 $-\frac12\tilde K\ln\frac{\zeta_z\zeta'_z}{\zeta_A\zeta_B}$，在 $\zeta_z\zeta'_z=Q^4$ 时为零；由\cref{thm:yzliu-intrinsic-soft}(a)，$S_r=S_I$，于是得\cref{eq:yzliu-dy-quasi}，$\hat\sigma=\mathcal H/(HH')$。在两个强子动量为 $\pm P^z$ 的质心系中，领头幂次下 $Q^2=xx's\simeq4xx'P_z^2$，故 $\zeta_z\zeta'_z=(2xP^z)^2(2x'P^z)^2=Q^4$ 自动成立。这里只用 $b$ 空间截面，不涉及 Fourier 符号约定。
\end{proof}

\paragraph{地位与缺口。}全阶论证 [A]，以\cref{thm:tmd-factorization}（[A]，含其 Glauber 论证）与\cref{thm:lamet-quasi-tmd}（[A]）为条件\cite[式~(22)]{Ji:2019sxk}\cite[式~(27)]{Ji:2019ewn}\cite[式~(155)]{Liu:2022nnk}。

\begin{proposition}[由准 TMD 之比得到 Sivers 不对称 \tagA]\label{prop:yzliu-sivers}
匹配核与自旋无关（只含硬胶子），故在同一 $(x,b_\perp,P^z)$ 处，准 Sivers 函数与非极化准 TMD 之比在幂次修正之内等于相应的光锥比值（$S_I$ 与 $H$ 相消）；DY 与 SIDIS 的符号来自订书钉方向 $\pm\hat z$，与\cref{thm:tmd-sivers-sign}一致。非微扰地位：\cref{asm:intro-bridge}、\cref{asm:intro-composite}（MA0、MA1、MA2）。
\uses{def:lamet-quasi-tmd,thm:lamet-quasi-tmd,prop:lamet-tmd-ratios,thm:tmd-sivers-sign}
\end{proposition}

\begin{proofsketch}
领头区域分析对一切自旋结构相同\cite[第~I--II~节]{Ji:2020jeb}：共线子图与光锥 TMD 相同，硬核只依赖色表示，故比值中 $H$、$S_I$ 与 CS 因子相消（\cref{prop:lamet-tmd-ratios}）。对沿 $+z$ 运动的强子，DY 型订书钉伸向 $+\hat z$、SIDIS 型伸向 $-\hat z$（\cref{def:lamet-quasi-tmd}），分别对应指向过去与指向未来的类光连线：指向过去的 $n_B$ 方向连线位于 $t<0$、$z>0$ 一侧，指向未来的位于 $t>0$、$z<0$ 一侧；由此继承\cref{thm:tmd-sivers-sign}的相对符号（这一相对符号与 $g$ 的符号约定无关）。比值在同一 $\vec b_\perp$ 处取，与 Fourier 符号约定无关。
\end{proofsketch}

\paragraph{地位与缺口。}全阶论证 [A]（概要）。证据：单圈、小 $b_\perp$ 极限下准 Sivers 函数再现共线扭度 3（Qiu--Sterman）描述；硬胶子极点（色因子 $\frac{C_A}2$）部分抵消软胶子极点（$\frac1{2N_c}$），端点与捏点项剩下 $\frac{C_A}2-\frac1{2N_c}=C_F$\cite[第~III~节]{Ji:2020jeb}。
